\section*{Appendix A: Proof of Proposition 1}
\label{app:prop1}
For any $i,j$ that $1\leq i,j\leq L$, we have
\begin{equation}
\begin{aligned}
    &\mathbb{E}[\eta_k(iT)\eta(jT)_k^*] \\
    &=\mathbb{E}\left[\int\mathbf{n}_c(t)\varphi^{*}(t-iT)dt\int\mathbf{n}_c^*(t)\varphi(t-jT)dt\right] \\
    &=\mathbb{E}\left[\int\int\mathbf{n}_c(t_1)\mathbf{n}_c(t_2)^{*}\varphi^{*}(t_1-iT)\varphi(t_2-jT)dt_1dt_2\right] \\
    &=\int\int\mathbb{E}[\mathbf{n}_c(t_1)\mathbf{n}_c(t_2)^{*}]\varphi^{*}(t_1-iT)\varphi(t_2-jT)dt_1dt_2 \\
    &=\int\int\sigma_C^2\delta(t_1-t_2)\varphi^{*}(t_1-iT)\varphi(t_2-jT)dt_1dt_2 \\
    &=\int\sigma_C^2\varphi^{*}(t-iT)\varphi(t-jT)dt \\
    &=\int\sigma_C^2\varphi^{*}(t)\varphi(t-(i-j)T)dt \\
    &=\sigma_C^2\phi((i-j)T)=\sigma_C^2\mathbf{\Phi}_{i,j}.
\end{aligned}
\end{equation}
Thus $\mathbb{E}[\bm{\eta}_k\bm{\eta}_k^H]=\sigma_C^2\mathbf{\Phi}$.

\section*{Appendix B: Proof of Proposition 2}
\label{app:prop2}
we can express $\overline{\mathbf{X}}_C$ as
\begin{equation}
    \overline{\mathbf{X}}_C = [\mathbf{E}_{1}\mathbf{X}_C^{\top},\mathbf{E}_{2}\mathbf{X}_C^{\top},\cdots,\mathbf{E}_{P}\mathbf{X}_C^{\top}],
\end{equation}
which yields a way to vectorize $\overline{\mathbf{X}}_C$ with respect to $\mathrm{vec}(\mathbf{S}^{\top})$
\begin{equation}
\begin{aligned}
    \mathrm{vec}(\overline{\mathbf{X}}_C) &= 
    \begin{bmatrix}
        \mathrm{vec}(\mathbf{E}_{1}\mathbf{X}_C^{\top}) \\
        \mathrm{vec}(\mathbf{E}_{2}\mathbf{X}_C^{\top}) \\
        \vdots \\
        \mathrm{vec}(\mathbf{E}_{P}\mathbf{X}_C^{\top})
    \end{bmatrix} = 
    \begin{bmatrix}
        \mathbf{I}_{N_t}\otimes\mathbf{E}_{1} \\
        \mathbf{I}_{N_t}\otimes\mathbf{E}_{2} \\
        \vdots \\
        \mathbf{I}_{N_t}\otimes\mathbf{E}_{P}
    \end{bmatrix}\mathrm{vec}(\mathbf{X}_C^{\top}) \\&= 
    \begin{bmatrix}
        \mathbf{I}_{N_t}\otimes(\mathbf{E}_{1}\mathbf{\Omega}_{\phi}) \\
        \mathbf{I}_{N_t}\otimes(\mathbf{E}_{2}\mathbf{\Omega}_{\phi}) \\
        \vdots \\
        \mathbf{I}_{N_t}\otimes(\mathbf{E}_{P}\mathbf{\Omega}_{\phi})
    \end{bmatrix}\mathrm{vec}(\mathbf{S}^{\top}).
\end{aligned}
\end{equation}
By using the fact that $\text{vec}(\mathbf{A}\mathbf{X}\mathbf{B})=(\mathbf{B}^{\top}\otimes\mathbf{A})\text{vec}(\mathbf{X})$, we have
\begin{equation}
\begin{aligned}
    \text{vec}(\mathbf{Y}_C^{\top}) = &(\mathbf{H}_C\otimes(\mathbf{G}\mathbf{U}_{\phi})^{\top})\text{vec}(\overline{\mathbf{X}}_C) + \text{vec}(\mathbf{N}_C^{\top}) \\
    = &\overline{\mathbf{H}}_C\text{vec}(\mathbf{S}^{\top}) + \text{vec}(\mathbf{N}_C^{\top}),
\end{aligned}
\end{equation}

\section*{Appendix C: Proof of Proposition 3}
\label{app:prop3}
Define $\mathbf{\Gamma}=\mathrm{Diag}([\Gamma_1,\Gamma_2,\cdots,\Gamma_K])$, $\overline{\mathbf{D}}=\mathrm{Diag}(\mathbf{\text{vec}(\mathbf{D}^{\top})})$ and $\bm{\varsigma}=\sqrt{\mathrm{diag}(\sigma_C^2\mathbf{I}_K\otimes\mathbf{\Lambda}_{\phi})}$; then, the CI constraint for $k$ users can be formulated as
\begin{equation}
\label{eq:mimo_ci}
\begin{aligned}
    \left| \Im\left\{\overline{\mathbf{D}}^{*}\overline{\mathbf{H}}_C\text{vec}(\mathbf{S}^{\top})\right\} \right| - \Re\left\{\overline{\mathbf{D}}^{*}\overline{\mathbf{H}}_C\text{vec}(\mathbf{S}^{\top})\right\}\tan\theta \\
    \leq (-\sqrt{\mathbf{\Gamma}\otimes\mathbf{I}_L}\tan\theta)\bm{\varsigma}.
\end{aligned}
\end{equation}
We first stack the $k$ linear inequalities
\begin{equation}
    \left| \Im\left\{\mathbf{d}_k^*\circ\mathbf{y}_{C,k}\right\} \right| - \Re\left\{\mathbf{d}_k^*\circ\mathbf{y}_{C,k}\right\}\tan\theta \leq (-\sqrt{\Gamma_k}\tan\theta)\bm{\sigma}
\end{equation}
in a column to form a united equality:
\begin{equation}
\label{eq:proof3_1}
\begin{aligned}
    &\left| \Im\left\{
    \begin{bmatrix}
        \mathbf{d}_1 \\
        \mathbf{d}_2 \\
        \vdots \\
        \mathbf{d}_k
    \end{bmatrix}^*
    \circ
    \begin{bmatrix}
        \mathbf{y}_{C,1} \\
        \mathbf{y}_{C,2} \\
        \vdots \\
        \mathbf{y}_{C,k}
    \end{bmatrix}\right\} \right| - 
    \Re\left\{
    \begin{bmatrix}
        \mathbf{d}_1 \\
        \mathbf{d}_2 \\
        \vdots \\
        \mathbf{d}_k
    \end{bmatrix}^*
    \circ
    \begin{bmatrix}
        \mathbf{y}_{C,1} \\
        \mathbf{y}_{C,2} \\
        \vdots \\
        \mathbf{y}_{C,k}
    \end{bmatrix}\right\}\tan\theta \\
    &=\left|\Im\left\{\text{vec}(\mathbf{D}^{\top})\circ\overline{\mathbf{H}}_C\text{vec}(\mathbf{S}^{\top}) \right\}\right| \\&- \Re\left\{\text{vec}(\mathbf{D}^{\top})\circ\overline{\mathbf{H}}_C\text{vec}(\mathbf{S}^{\top}) \right\}\tan\theta \leq
    \begin{bmatrix}
        (-\sqrt{\Gamma_1}\tan\theta)\bm{\varsigma} \\
        (-\sqrt{\Gamma_2}\tan\theta)\bm{\varsigma} \\
        \vdots \\
        (-\sqrt{\Gamma_k}\tan\theta)\bm{\varsigma}
    \end{bmatrix}.
\end{aligned}
\end{equation}
Notice that for Hardmard product of two vectors we have $\mathbf{a}\circ\mathbf{b}=\text{Diag}(\mathbf{a})\mathbf{b}$. Thus we can rewrite (\ref{eq:proof3_1}) as
\begin{equation}
\begin{aligned}
    &\left|\Im\left\{\text{Diag}(\text{vec}(\mathbf{D}^{\top}))\overline{\mathbf{H}}_C\text{vec}(\mathbf{S}^{\top}) \right\}\right| \\ &\qquad\qquad- \Re\left\{\text{Diag}(\text{vec}(\mathbf{D}^{\top}))\overline{\mathbf{H}}_C\text{vec}(\mathbf{S}^{\top}) \right\}\tan\theta \\
    &=\left| \Im\left\{\overline{\mathbf{D}}^{*}\overline{\mathbf{H}}_C\text{vec}(\mathbf{S}^{\top})\right\} \right| - \Re\left\{\overline{\mathbf{D}}^{*}\overline{\mathbf{H}}_C\text{vec}(\mathbf{S}^{\top})\right\}\tan\theta \\
    &\leq
    \begin{bmatrix}
        (-\sqrt{\Gamma_1}\tan\theta)\bm{\varsigma} \\
        (-\sqrt{\Gamma_2}\tan\theta)\bm{\varsigma} \\
        \vdots \\
        (-\sqrt{\Gamma_k}\tan\theta)\bm{\varsigma}
    \end{bmatrix}
    =(-\sqrt{\mathbf{\Gamma}\otimes\mathbf{I}_L}\tan\theta)\bm{\varsigma}.
\end{aligned}
\end{equation}


\section*{Appendix D: Proof of Proposition 4}
\label{app:prop4}
\begin{equation}
\begin{aligned}
    &\int\Vert x_n(t)\Vert^2 dt = \int\left\Vert\sum_{i=1}^{L}\varphi(t-(i-1)T)s_{n,i}\right\Vert^2 dt \\
    &=\int\left(\sum_{i=0}^{L-1}\varphi(t-iT)s_{n,i+1}\right)\left(\sum_{i=0}^{L-1}\varphi(t-iT)^*s_{n,i+1}^*\right) dt \\
    &=\int\sum_{i=1}^{L}\sum_{j=1}^{L}\left(\varphi(t-(i-1)T)\varphi(t-(j-1)T)^*s_{n,i}s_{n,j}^*\right) dt \\
    &=\sum_{i=1}^{L}\sum_{j=1}^{L}\left(\int\varphi(t-(i-1)T)\varphi(t-(j-1)T)^*dt\right)s_{n,i}s_{n,j}^* dt\\
    &=\sum_{i=1}^{L}\sum_{j=1}^{L}\left(\int\varphi(t)\varphi(t-(i-j)T)^*dt\right)s_{n,i}s_{n,j}^* dt\\
    &=\sum_{i=1}^{L}\sum_{j=1}^{L}\phi((i-j)T)s_{n,i}s_{n,j}^* = \mathbf{s}_n^{H}\mathbf{\Phi}\mathbf{s}_n.
\end{aligned}
\end{equation}

\section*{Appendix E: Proof of Proposition 5}
\label{app:prop5}
Similar to the vectorization of $\overline{\mathbf{X}}_C$, we are able to vectorize $\overline{\mathbf{X}}$ as
\begin{equation}
\begin{aligned}
    \text{vec}(\overline{\mathbf{X}}_R) &= \text{vec}([\mathbf{E}_1\mathbf{X}_R^{\top},\mathbf{E}_2\mathbf{X}_R^{\top},\cdots,\mathbf{E}_P\mathbf{X}_R^{\top}])\\
    &= \begin{bmatrix}
        \mathbf{I}_N\otimes(\mathbf{E}_{1}\mathbf{\Omega}_{\varphi}) \\
        \mathbf{I}_N\otimes(\mathbf{E}_{2}\mathbf{\Omega}_{\varphi}) \\
        \vdots \\
        \mathbf{I}_N\otimes(\mathbf{E}_{P}\mathbf{\Omega}_{\varphi})
    \end{bmatrix}\text{vec}(\mathbf{S}^{\top}) = \mathbf{E}_R\text{vec}(\mathbf{S}^{\top}).
\end{aligned}
\end{equation}
Using the fact that $\text{tr}(\mathbf{A}^H\mathbf{B})=\text{vec}^H(\mathbf{A})\text{vec}(\mathbf{B})$ and $\text{tr}(\mathbf{ABCD})=\text{vec}^{\top}(\mathbf{D})(\mathbf{A}\otimes\mathbf{C}^{\top})\text{vec}(\mathbf{B}^{\top})$, we have
\begin{equation}
\begin{aligned}
    \text{tr}(\mathbf{Q}_k^H\overline{\mathbf{X}}_R) &= \text{vec}^H(\mathbf{Q}_k)\text{vec}(\overline{\mathbf{X}}_R) \\
    &= \text{vec}^H(\mathbf{Q}_k)\mathbf{E}_R\text{vec}(\mathbf{S}^{\top})
    = -\mathbf{b}_k^H\text{vec}(\mathbf{S}^{\top}),
\end{aligned}
\end{equation}
\begin{equation}
\begin{aligned}
    \text{tr}(\mathbf{T}_k\overline{\mathbf{X}_R}\overline{\mathbf{X}}_R^H) &= \text{tr}(\mathbf{I}_L\overline{\mathbf{X}}_R^H\mathbf{T}_k\overline{\mathbf{X}}_R) \\
    &= \text{vec}^{\top}(\overline{\mathbf{X}}_R)(\mathbf{I}_L\otimes\mathbf{T}_k^*)\text{vec}(\overline{\mathbf{X}}_R^*) \\
    &= \text{vec}(\mathbf{S}^{\top})^{\top}\mathbf{E}_R^{\top}(\mathbf{I}_L\otimes\mathbf{T}_k^*)\mathbf{E}_R^*\text{vec}(\mathbf{S}^{\top})^* \\
    &= \text{vec}(\mathbf{S}^{\top})^H\mathbf{E}_R^H(\mathbf{I}_L\otimes\mathbf{T}_k)\mathbf{E}_R\text{vec}(\mathbf{S}^{\top}) \\
    &= \text{vec}(\mathbf{S}^{\top})^H\mathbf{B}_k\text{vec}(\mathbf{S}^{\top})/\sigma_H^2.
\end{aligned}
\end{equation}
Thus we have
\begin{equation}
\begin{aligned}
    &2\Re\left\{\text{tr}(\mathbf{Q}_k^H\overline{\mathbf{X}}_R)\right\} - \text{tr}(\mathbf{T}_k(\sigma_H^2\overline{\mathbf{X}}_R\overline{\mathbf{X}}_R^H+\sigma_R^2\mathbf{I})) \\
    &= \text{tr}(\sigma_R^2\mathbf{T}_k) + 2\Re\left\{\text{tr}(\mathbf{Q}_k^H\overline{\mathbf{X}}_R)\right\} - \sigma_H^2\text{tr}(\mathbf{T}_k\overline{\mathbf{X}}_R\overline{\mathbf{X}}_R^H) \\
    &= c_k - 2\Re\left\{\text{vec}(\mathbf{S}^{\top})^H\mathbf{b}_k\right\}-\text{vec}(\mathbf{S}^{\top})^H\mathbf{B}_k\text{vec}(\mathbf{S}^{\top}).
\end{aligned}
\end{equation}

\section*{Appendix F: Proof of Proposition 6}
\label{app:prop6}
The proof of the convergence of the BPS algorithm to the optimal solution is presented below. Consider the following convex problem:
\begin{equation}
\label{opt:prov_origin}
\begin{aligned}
    &\underset{\mathbf{x}}{\min} \;
    f_o(\mathbf{x}) \\
    &s.t.\; \mathbf{x}\in\mathcal{S},\; f_p(\mathbf{x})\leq\mathcal{E}.
\end{aligned}
\end{equation}
where $f_o$ and $f_p$ are convex functions and $\mathcal{S}$ is a convex region.
Also consider the penalty problem, which eliminates the energy constraint
\begin{equation}
\label{opt:prov_penalty}
\begin{aligned}
    \mathcal{P}(\rho):\; &\underset{\mathbf{x}}{\min} \;
    f_o(\mathbf{x})+\rho f_p(\mathbf{x}) \\
    &s.t.\; \mathbf{x}\in\mathcal{S}.
\end{aligned}
\end{equation}
Initially, we will prove that as $\rho$ increases, the optimal solution of $\mathcal{P}(\rho)$, $\mathbf{x}^{\star}$, exhibits a larger $f_o(\mathbf{x}^{\star})$ and a smaller $f_p(\mathbf{x}^{\star})$. Subsequently, we will establish that the optimal solution of problem (\ref{opt:prov_origin}) is also the optimal solution of problem $\mathcal{P}(\rho)$ for a certain $\rho$.
\begin{align}
    f_o(\mathbf{x}_1^{\star})+\rho_2f_p(\mathbf{x}_1^{\star})\geq f_o(\mathbf{x}_2^{\star})+\rho_2f_p(\mathbf{x}_2^{\star}), \label{eq:prov1}\\
    f_o(\mathbf{x}_1^{\star})+\rho_1f_p(\mathbf{x}_1^{\star})\leq f_o(\mathbf{x}_2^{\star})+\rho_1f_p(\mathbf{x}_2^{\star}). \label{eq:prov2}
\end{align}
Then we rearrange the inequality (\ref{eq:prov2}) as
\begin{equation}
\label{eq:prov3}
\begin{aligned}
    f_o(\mathbf{x}_1^{\star})+\rho_2f_p(\mathbf{x}_1^{\star})+(\rho_1-\rho_2)f_p(\mathbf{x}_1^{\star}) \\ \leq f_o(\mathbf{x}_2^{\star})+\rho_2f_p(\mathbf{x}_2^{\star})+(\rho_1-\rho_2)f_p(\mathbf{x}_2^{\star}).
\end{aligned}
\end{equation}
According to inequality (\ref{eq:prov1}) and (\ref{eq:prov3}) we have
\begin{equation}
\begin{aligned}
    0\leq\left(f_o(\mathbf{x}_1^{\star})+\rho_2f_p(\mathbf{x}_1^{\star})\right) - \left(f_o(\mathbf{x}_2^{\star})+\rho_2f_p(\mathbf{x}_2^{\star})\right)
    \\ \leq
    (\rho_1-\rho_2)(f_p(\mathbf{x}_2^{\star})-f_p(\mathbf{x}_1^{\star})),
\end{aligned}
\end{equation}
which yields
\begin{equation}
\label{eq:prov4}
    f_p(\mathbf{x}_1^{\star})\leq f_p(\mathbf{x}_2^{\star}).
\end{equation}
Then combining inequality (\ref{eq:prov1}) and (\ref{eq:prov4}) we have
\begin{equation}
\label{eq:prov5}
    f_o(\mathbf{x}_1^{\star})\geq f_o(\mathbf{x}_2^{\star})
\end{equation}
Inequality (\ref{eq:prov4}) and (\ref{eq:prov5}) reveal that an increase in the penalty factor results in an increment of the objective function and a reduction in the penalty function. The objective of the BPS algorithm is to identify the minimal penalty factor such that the penalty function $f_p$ does not surpass $\mathcal{E}$.

We now prove that the optimal solution $\mathbf{x}^{\star}$ of problem (\ref{opt:prov_origin}) corresponds to the optimal solution of problem $\mathcal{P}(\rho)$ for a particular $\rho$. Assuming $\mathbf{x}^{\star}$ is not constrained by $f_p(\mathbf{x})\leq\mathcal{E}$, it is evident that $\mathbf{x}^{\star}$ represents the optimal solution of $\mathcal{P}(0)$. If it is constrained, in accordance with the KKT complementary slackness condition, we ascertain that for certain values of $\mu$ and $\bm{\nu}$:
\begin{equation}
    \frac{\partial f_o}{\partial\mathbf{x}}(\mathbf{x}^{\star})+\mu\frac{\partial f_p}{\partial\mathbf{x}}(\mathbf{x}^{\star})+\frac{\partial{\mathcal{L}_{\mathcal{S}}(\bm{\nu})}}{\partial\mathbf{x}}(\mathbf{x}^{\star})=0,
\end{equation}
where $\mathcal{L}_{\mathcal{S}}$ and $\bm{\nu}$ denote the Lagrange augmentation function and dual variables for the constraint $\mathbf{x}\in\mathcal{S}$. Consequently, we deduce that this condition also represents the KKT complementary slackness condition for problem $\mathcal{P}(\mu)$. Therefore, $\mathbf{x}^{\star}$ serves as the optimal solution of $\mathcal{P}(\mu)$.